\documentclass[12pt]{ctexart}
\usepackage[a4paper, margin=2.5cm]{geometry}
\usepackage{array}
\usepackage{longtable}
\usepackage{booktabs}
\usepackage{multirow}
\usepackage{tcolorbox}
\usepackage{enumitem}
\usepackage{fancyhdr}
\usepackage{lastpage}

% 设置字体
\setmainfont{SimSun}
\setCJKmainfont{SimSun}

% 自定义命令定义
\newcommand{\projecttitle}[1]{\def\@projecttitle{#1}}
\newcommand{\groupname}[1]{\def\@groupname{#1}}
\newcommand{\groupmembers}[1]{\def\@groupmembers{#1}}
\newcommand{\instructor}[1]{\def\@instructor{#1}}
\newcommand{\dateperiod}[1]{\def\@dateperiod{#1}}

% 默认值
\projecttitle{\texttt{[填入项目名称，例如：校园二手交易微信小程序]}}
\groupname{\texttt{[例如：xxxx]}}
\groupmembers{\texttt{[小组成员姓名列表]}}
\instructor{\texttt{[教师姓名]}}
\dateperiod{\texttt{202X年X月X日 至 202X年X月X日}}

% 页脚设置
\pagestyle{fancy}
\fancyhf{}
\rfoot{\today\ | 模板版本: V1.0}
\cfoot{\thepage/\pageref{LastPage}}

\title{\@projecttitle}
\author{}
\date{}

\begin{document}

% 封面
\begin{titlepage}
    \centering
    \vspace*{2cm}
    {\LARGE\bfseries 实验一 项目计划书\par}
    \vspace{3cm}
    \begin{tabular}{ll}
        项目名称： & \@projecttitle \\
        小组名称： & \@groupname \\
        小组成员： & \@groupmembers \\
        指导教师： & \@instructor \\
        起止日期： & \@dateperiod \\
    \end{tabular}
    \vfill
    \today
\end{titlepage}

\newpage

% 版本修订记录
\section*{版本修订记录}
\begin{center}
\begin{tabular}{|p{2cm}|p{2cm}|p{4cm}|p{2cm}|p{2cm}|}
\hline
\textbf{版本号} & \textbf{修改日期} & \textbf{修改内容} & \textbf{修改人} & \textbf{审核人} \\
\hline
V1.0 & 202X.XX.XX & 初稿创建，明确需求与分工 & 张三 & 李四 \\
\hline
\texttt{[...]} & \texttt{[...]} & \texttt{[...]} & \texttt{[...]} & \texttt{[...]} \\
\hline
\end{tabular}
\end{center}
% 可动态添加更多修订行

\newpage

% 1. 引言
\section{引言}

% 1.1 项目背景与目标
\subsection{项目背景与目标}
% 撰写要求：简明扼要阐述项目动因与解决痛点
\begin{quote}
目前校内二手交易多依赖 QQ 群和贴吧，信息杂乱、无分类、交易无保障。本项目旨在开发一款针对本校的实名制二手交易小程序，为学生提供安全、便捷的闲置物品流转平台。
\end{quote}

\texttt{[在此处撰写项目背景与目标...]}

% 1.2 项目范围
\subsection{项目范围}
% 撰写要求：清晰界定功能边界，包括核心功能、非功能需求和排除项
\begin{itemize}
    \item \textbf{核心功能：}
    \begin{itemize}
        \item 用户微信授权登录
        \item 商品发布（图文）
        \item 分类浏览
        \item 私信聊天
        \item 我的发布管理
        \item \texttt{[在此处添加其他核心功能...]}
    \end{itemize}
    
    \item \textbf{非功能范围：}
    \begin{itemize}
        \item 支撑至少 1000 人同时使用
        \item 页面加载速度小于 2 秒
        \item \texttt{[在此处添加其他非功能需求...]}
    \end{itemize}
    
    \item \textbf{排除项：}
    \begin{itemize}
        \item 暂不涉及在线支付功能，仅做线下交易引导
        \item 暂不做后台管理员端，仅做普通用户端
        \item \texttt{[在此处添加其他排除项...]}
    \end{itemize}
\end{itemize}

% 1.3 术语与缩写
\subsection{术语与缩写}
% 撰写要求：解释项目中使用的专业术语和缩写
\begin{description}
    \item[IM] Instant Messaging，即时通讯
    \item[UI] User Interface，用户界面
    \item[API] Application Programming Interface，应用程序接口
    \item[\texttt{[...]}] \texttt{[在此处添加其他术语...]}
\end{description}

% 2. 项目组织结构
\section{项目组织结构}
% 撰写要求：明确团队成员的角色分工和联系方式
\begin{center}
\begin{tabular}{|p{2cm}|p{2cm}|p{2.5cm}|p{3.5cm}|p{2cm}|}
\hline
\textbf{姓名} & \textbf{学号} & \textbf{角色} & \textbf{主要职责} & \textbf{联系方式} \\
\hline
张三 & 2021001 & 项目经理/组长 & 进度控制、需求分析、文档管理、UI设计 & 138xxxx \\
\hline
李四 & 2021002 & 后端开发工程师 & 数据库设计、API接口开发、服务器部署 & 139xxxx \\
\hline
王五 & 2021003 & 前端开发工程师 & 小程序页面实现、API对接、用户体验优化 & 137xxxx \\
\hline
\texttt{[...]} & \texttt{[...]} & \texttt{[...]} & \texttt{[...]} & \texttt{[...]} \\
\hline
\end{tabular}
\end{center}

\begin{tcolorbox}
\textbf{注意：} 角色可以有兼任，但必须职责分明
\end{tcolorbox}

% 3. 技术与过程规范
\section{技术与过程规范}

% 3.1 技术选型
\subsection{技术选型}
% 撰写要求：说明各技术层的选型理由和具体技术栈
\begin{itemize}
    \item \textbf{前端：} \texttt{[例如：微信小程序原生框架 / Vue3 + Element Plus / React Native]}
    \item \textbf{后端：} \texttt{[例如：Python + Django / Java + Spring Boot / Node.js + Express]}
    \item \textbf{数据库：} \texttt{[例如：MySQL (关系型) / MongoDB (非关系型) / Redis (缓存)]}
    \item \textbf{服务器与部署：} \texttt{[例如：腾讯云/阿里云轻量服务器 + Docker + Nginx]}
\end{itemize}

% 3.2 开发规范
\subsection{开发规范}
% 撰写要求：制定团队协作和代码质量的标准规范
\begin{enumerate}
    \item \textbf{版本控制：} 指定 Git 仓库地址（如 GitHub/Gitee），约定分支管理策略（例如：main为主分支，dev为开发分支，feature-xxx为功能分支）
    \item \textbf{编码规范：} 参照 [PEP8 (Python) / 阿里Java开发手册] 等，确保代码可读性
    \item \textbf{文档规范：} 所有接口需编写 API 文档（推荐使用 Swagger 或 Postman 文档功能）
\end{enumerate}

% 4. 项目进度计划
\section{项目进度计划}

% 4.1 里程碑计划
\subsection{里程碑计划}
% 撰写要求：规划关键节点和交付产物
\begin{center}
\begin{tabular}{|p{4cm}|p{3cm}|p{4cm}|}
\hline
\textbf{里程碑} & \textbf{预计完成日期} & \textbf{交付产物} \\
\hline
需求分析完成 & 第3周 & 需求规格说明书、原型图 \\
\hline
系统设计完成 & 第5周 & 数据库E-R图、架构设计文档、API设计文档 \\
\hline
核心功能开发完成（Alpha版） & 第9周 & 可运行的软件包、单元测试报告 \\
\hline
测试与Bug修复完成（Beta版） & 第11周 & 测试报告、用户手册 \\
\hline
最终交付与演示 & 第12周 & 项目源码、部署文档、答辩PPT \\
\hline
\end{tabular}
\end{center}

% 4.2 详细任务分解（WBS - 迭代1：基础功能为例）
\subsection{详细任务分解（WBS - 迭代1：基础功能为例）}
% 撰写要求：细化到具体任务，明确负责人、工时和依赖关系
\begin{center}
\begin{tabular}{|p{1.5cm}|p{3cm}|p{2cm}|p{1.5cm}|p{2cm}|p{2cm}|p{2cm}|}
\hline
\textbf{任务ID} & \textbf{任务名称} & \textbf{负责人} & \textbf{预估工时} & \textbf{开始日期} & \textbf{结束日期} & \textbf{前置任务} \\
\hline
1.1 & 数据库表结构创建（用户表、商品表） & 李四 & 4h & 第5周一 & 第5周一 & 无 \\
\hline
1.2 & 注册/登录接口开发 & 李四 & 6h & 第5周二 & 第5周三 & 1.1 \\
\hline
1.3 & 登录页面UI开发 & 王五 & 4h & 第5周二 & 第5周二 & 无 \\
\hline
1.4 & 登录功能前后端联调 & 李四、王五 & 2h & 第5周四 & 第5周四 & 1.2, 1.3 \\
\hline
\texttt{[...]} & \texttt{[...]} & \texttt{[...]} & \texttt{[...]} & \texttt{[...]} & \texttt{[...]} & \texttt{[...]} \\
\hline
\end{tabular}
\end{center}
% 如需其他迭代，请复制此表格并修改标题

% 5. 风险管理计划
\section{风险管理计划}
% 撰写要求：识别潜在风险并制定应对措施
\begin{center}
\begin{tabular}{|p{4cm}|p{2cm}|p{2cm}|p{5cm}|}
\hline
\textbf{风险描述} & \textbf{可能性 (高/中/低)} & \textbf{影响程度 (高/中/低)} & \textbf{应对措施} \\
\hline
技术风险： 小组成员未使用过微信小程序开发，学习成本高导致进度延误。 & 高 & 高 & 
\begin{enumerate}
    \item 提前一周开始技术预研。
    \item 组内组织技术分享会。
\end{enumerate} \\
\hline
人员风险： 考试周组员忙于考试，投入项目时间不足。 & 中 & 高 & 
\begin{enumerate}
    \item 制定详细分工，避免任务耦合过紧。
    \item 提前预留缓冲时间。
\end{enumerate} \\
\hline
需求风险： 用户（同学）反馈需求不断变更。 & 中 & 中 & 
\begin{enumerate}
    \item 严格遵循变更控制流程。
    \item 第一期必须锁定范围，新增功能放入V2.0计划。
\end{enumerate} \\
\hline
接口风险： 依赖的第三方API（如地图）服务不稳定。 & 低 & 中 & 
\begin{enumerate}
    \item 调研备选API方案。
    \item 增加接口调用失败的重试和降级处理逻辑。
\end{enumerate} \\
\hline
\texttt{[...]} & \texttt{[...]} & \texttt{[...]} & \texttt{[在此处添加其他风险及应对措施...]} \\
\hline
\end{tabular}
\end{center}

% 6. 资源配置与支持
\section{资源配置与支持}
% 撰写要求：列出所需的外部资源和支持
\begin{tcolorbox}
\textbf{示例：} 需要申请一台云服务器用于部署后端；需要使用实验室的服务器进行压力测试；需要向学院申请测试账号数据
\end{tcolorbox}

\texttt{[在此处列出所需外部资源...]}

% 7. 审批意见
\section{审批意见}
% 撰写要求：提供签字确认区域
\begin{center}
\begin{tabular}{|p{3cm}|p{3cm}|p{3cm}|p{3cm}|}
\hline
\textbf{角色} & \textbf{姓名} & \textbf{签字} & \textbf{日期} \\
\hline
项目组长 & \texttt{[...]} & \texttt{[...]} & \texttt{[...]} \\
\hline
指导教师 & \texttt{[...]} & \texttt{[...]} & \texttt{[...]} \\
\hline
\end{tabular}
\end{center}

% 附录：评分要点分析（仅作为参考，不渲染在正文中）
% \appendix
% \section*{附录：评分要点分析}
% \begin{center}
% \begin{tabular}{|p{2cm}|p{3cm}|p{2cm}|p{8cm}|}
% \hline
% \textbf{编号} & \textbf{评分项} & \textbf{权重} & \textbf{评价标准说明} \\
% \hline
% 1 & 可行性 & 30\% & 技术选型是否合理？任务量是否能在课程周期内完成？避免题目过难导致烂尾，或过简单缺乏工作量 \\
% \hline
% 2 & 明确性 & 25\% & 分工是否清晰无歧义？项目范围（边界）是否界定清楚？ \\
% \hline
% 3 & 可追踪性 & 25\% & 时间节点是否具体？WBS任务分解是否细化到了“小时”或“人天”？任务粒度过大则无法有效跟踪进度 \\
% \hline
% 4 & 风险意识 & 20\% & 是否识别出了团队真实可能遇到的困难？是否有备选方案？这是软件工程区别于“写代码”的重要标志 \\
% \hline
% \end{tabular}
% \end{center}

\end{document}
